\clearpage\fancyhead[L]{\small REVIEWS. LOGARITHM TABLES.}
\SourceLeaf{262}
\pdfbookmark[1]{Babbage. 1828 January 19.}{babbage1828}\par\Needspace{10\baselineskip}\begin{center}\rule{0.7\linewidth}{0.4pt}\par Göttingische gelehrte Anzeigen. 1828 January 19.\par\rule{0.7\linewidth}{0.4pt}\end{center}

\noindent\textit{Table of logarithms of the natural numbers, from} 1 \textit{to} 108000, \textit{by} \textsc{Charles}
\textsc{Babbage.} \textit{Stereotyped. London. Printed for J. Mawman.} 1827. 202 pp., large 8vo.

This new printing of the logarithm tables is distinguished from others by more deliberate attention to small incidental details. Those who only occasionally need to look up a few logarithms in the tables mainly ask for the greatest possible accuracy. But for others, to whom the tables are an everyday working tool, even the slightest details that can affect convenience of use cease to be matters of indifference. The colour, thickness and beauty of the paper; the format; the size, sharpness and attractive shape of the type; the quality of the printing ink; the arrangement of the numbers so that what is wanted can be found quickly and reliably without tiring the eye; the presence of everything needed, but also the absence of everything one does not want and that would merely make the table harder to take in: all these details acquire some importance in a task repeated a hundred times a day. Admittedly, much depends on the individual calculator and their habits, so that no single edition can well suit everyone best. A large format is cumbersome for a nearsighted person, who prefers smaller type, while for a farsighted person the reverse is true. A less practised calculator values one addition or another that a more practised calculator, who performs several parts of the operations mentally without effort, would rather have omitted, since everything superfluous is merely distracting. Yet there are also rules of general validity. In the preface, the editor of the present tables has collected a dozen such rules, derived from comparing many logarithm tables. Most are immediately self-evident, although some may deserve acceptance only with qualifications.

What most distinguishes the present edition of the logarithm tables from others is that it is printed on coloured (yellow) paper: one soon becomes accustomed to it and finds it, especially for use by artificial light, more \SourceLeaf{263} pleasant than white paper. The arrangement is essentially the usual one. After a change in the third digit, the fourth digit in the remaining logarithms in the same row is set in different type, about half the size of the rest. For the reviewer, who generally attaches little value to such warning signs, this type has not lost its unfamiliar, distracting appearance even after a couple of months' use. He would prefer the asterisks used in \textsc{Vega}’s table, or the dots printed in the better copies of \textsc{Taylor}’s tables. The editor chose not to use asterisks because they would have made the columns too wide, leaving no room for converting the numbers, regarded as seconds, into degrees, minutes and seconds. The reviewer would gladly have dispensed with this conversion column (arranged just as in \textsc{Callet}’s tables). A practised calculator will not open the tables for a conversion so easily made mentally; it is useful, however, that the editor has at least separated these columns from the logarithm table by a heavy line. Finally, where the last digit has been increased (because the logarithm, if continued, would have had 5 or a larger number as its eighth digit), this is marked by a dot placed beneath that seventh digit. This arrangement is at least better than \textsc{Prasse}'s use of differently shaped type; yet even so, some digits retain an unpleasantly unfamiliar appearance which, in our view already expressed in the review of \textsc{Prasse}’s tables in these pages [1811 May 25], is not outweighed by the benefit.

Otherwise, the type and paper in this edition of the logarithm tables leave nothing to be desired, and the greatest care has been devoted to accuracy. In an edition that gives so much care to small details affecting convenience of use, the sparse provision of proportional parts is surprising: one often has to turn backwards or forwards to find those needed at the moment.
\SourceLeaf{264}
